% ============================================================
% Section 6: Conclusion
% ============================================================

\section{Conclusion}
\label{sec:conclusion}

We presented DocVerify, a multi-task deep learning system for joint detection and
localization of forged identity documents, and a multi-objective Pareto HPO
framework that selects hyperparameter configurations by optimizing detection
quality (PR-AUC) and localization quality (Dice) simultaneously, without fixing
their relative importance in advance.

Our main findings are:

\begin{enumerate}
  \item Joint multitask training with Pareto-optimized loss weights achieves
  significantly higher Dice than both single-task variants and equal-weight
  joint training (blind test, 30 seeds, $p = 1.3 \times 10^{-4}$, $d = 1.01$),
  while maintaining competitive PR-AUC.

  \item Under a rigorous 10-fold nested cross-validation protocol with 50 MOTPE
  trials per fold, DocVerify achieves PR-AUC $= 0.9921 \pm 0.0058$ and
  Dice $= 0.856 \pm 0.030$ on held-out folds, demonstrating consistent
  generalization across data partitions.

  \item Training exclusively on FantasyID and without any external pre-training,
  DocVerify achieves results competitive with the top-3 entries of the
  DeepID 2025 Challenge, surpassing the challenge winner on Track~2
  localization ($0.807$ vs.\ $0.784$) and all FantasyID-only systems on
  Track~1 detection.

  \item Multi-objective Pareto HPO is statistically indistinguishable from
  scalar selection on absolute metrics ($n = 10$, all $p \geq 0.05$), but
  provides principled, objective-agnostic model selection and a richer
  characterization of the HPO objective landscape.
\end{enumerate}

The source code, training scripts, and all experimental artifacts are publicly
available at \url{https://github.com/HernanDiaz/IDVerify/}.

Several directions remain open. The most pressing is improving
cross-domain generalization: models trained on a single synthetic
dataset cannot be assumed to transfer automatically to real-world document
designs, and evaluating robustness across domains remains future work.
Three complementary strategies are worth investigating:
(i)~\emph{forgery-artifact-centric features} --- noise fingerprints,
compression inconsistencies, and frequency-domain anomalies~\cite{noiseprint2020}
that are independent of document appearance and thus naturally domain-invariant;
(ii)~\emph{foundation model backbones} such as DINOv2 or CLIP, whose
large-scale pre-training provides richer general-purpose representations
that have already shown cross-domain benefits in concurrent
work~\cite{fakeidet2025};
(iii)~\emph{domain-adversarial training}~\cite{ganin2016dann}, where a
domain discriminator forces the encoder to learn features that are
invariant across document types.
An official submission to the DeepID 2025 Challenge test server would also
provide a directly controlled comparison with published baselines,
eliminating the internal-holdout caveat noted in Section~\ref{sec:discussion}.

% ── Acknowledgments ──────────────────────────────────────────
\section*{Acknowledgments}
This work was supported by the Instituto Nacional de Ciberseguridad
(INCIBE, Spain's National Cybersecurity Institute) and the INCIBE Chair
in Cybersecurity and Artificial Intelligence at the University of Oviedo, Spain. The author thanks Rafael Gonz\'{a}lez Quesada for his
participation in the project, and the Idiap Research Institute for releasing the
FantasyID dataset and the DeepID 2025 Challenge organizers for providing the
evaluation framework.
